\section{Weighted tensor cycles with an entire external partner}
\label{sec:tensor-weighted-partner-extension}

\paragraph{Checked status and exact interface.}
HD-TW04--05, question Q2du, is a new root derivation extending the accepted local weighted-cycle argument. The author and independent reviewer have read the complete proof and passed it without repair. A separate exact weighted-orbit and polynomial supplement passed; the universal proof does not rely on finite sampling. An actual partner support
hypothesis is retained throughout; it is not inferred from an incomplete
factor catalogue.

Let an actual finite-group sphere cover have a genuine generating
product-one inertia tuple \(T_1,\ldots,T_r\), coefficient space
\[
 V=W^{\otimes t}\otimes U,\qquad
 a=\dim W\ge4,\quad b=\dim U\ge1,\quad e=a^t b,
 \qquad r\ge12,\quad \sum_i\operatorname{rank}(1-T_i)=2e+6.
 \tag{TWP.1}
\]
Here \(t\ge2\). The factor \(W^{\otimes t}\) is stable as a whole,
and its actual projective action normalizes the ordinary core \(A_n^t\)
and permutes its factors transitively. The compatible cycle products
are represented on \(W\) by an irreducible ordinary \(S_n\)-module,
up to scalars. Both the action on the \(t\) factors and the natural
action on \(nt\) letters are actual quotient actions of this cover.
The action on \(U\) is an actual projective factor action.

For a semisimple operator write \(\nu(S)=\dim S-\max_z\dim\ker(S-z)\).
Assume that every nonscalar operator of the actual partner image satisfies
\[
                         \nu(S)/b\ge\eta>0.
                         \tag{TWP.2}
\]
If \(b=1\), this condition is vacuous; only genuine total scalar
entries must be considered below.

\subsection*{HD-TW04. The complete weighted-cycle extension}

Suppose every nonidentity \(x\in S_n\) satisfies
\[
                    \nu(W(x))\ge A\operatorname{ind}(x)+B,
                    \tag{TWP.3}
\]
where
\[
 \begin{gathered}
 A\ge0,\quad B>0,\quad 2B\le a-1,\quad
 \eta\ge B/a,\\
 \lambda=1-(nA+1)/a\ge0,\qquad L=a-(n-1)A\ge1,\\
                 a(8B-2)>6,\qquad a(9B-L-1)>6.
 \end{gathered}\tag{TWP.4}
\]
Then the source (TWP.1) does not exist.

\begin{proof}
Tensor spectral convolution bounds the normalized moving rank of a total
entry below by the normalized projective support of its operator on
\(W^{\otimes t}\), and also by that on \(U\). Every scalar phase is
allowed in this assertion.

Consider a factor-fixing entry. If at least one of its within-factor
permutations is nonidentity, average (TWP.3) over its \(m\ge1\) active
factors. If \(I_i\) is the sum of their natural indices, the maximum
of their normalized supports is at least
\(AI_i/(ma)+B/a\ge AI_i/(ta)+B/a\).
If all these permutations are identity, the whole first factor is scalar
and \(I_i=0\). A nonscalar partner then costs at least \(\eta\ge B/a\).
If the partner is scalar too, genuineness means that the total operator
is a nonidentity scalar, with full moving rank. Since \(B/a<1\), that
case has the same lower bound. Hence every genuine factor-fixing entry
satisfies
\[
       \operatorname{rank}(1-T_i)/e\ge AI_i/(ta)+B/a.
       \tag{TWP.5}
\]
This is the step that was unavailable for an unrestricted external partner.

A factor-moving entry has a tensor cycle of length at least two. After
diagonalizing its finite-projective-order cycle product, constant words
number \(a\), and every other word orbit has length at least two.
Every eigenspace on such a weighted orbit has dimension at most one.
Thus that entry has normalized support at least \((a-1)/(2a)\).
Convolution with all remaining cycles and the external partner preserves
this bound. In particular every genuine entry costs at least \(B/a\),
by \(2B\le a-1\).

If four entries move factors, choose four such entries and eight of the
remaining genuine entries. Their total moving rank is at least
\[
       4\frac{(a-1)e}{2a}+8\frac{Be}{a}
         =2e+(8B-2)e/a
         \ge2e+a(8B-2)>2e+6,
         \tag{TWP.6}
\]
because \(t\ge2\) and \(b\ge1\). Thus at most three entries move
factors, and at least nine fix them.

For an arbitrary entry, let \(\pi_i\) be its factor permutation.
Its cycles have lengths \(\ell_j\) and represented products \(x_j\).
The accepted exact weighted-cycle bound of D-ATT08, with its proved
short-word estimate, gives on that cycle
\[
 s_j\ge1-1/\ell_j+
       \nu(W(x_j))/(\ell_j a)
       -R_{\ell_j}(a)/(\ell_j a^{\ell_j}),
 \quad R_1=0,\quad
 R_\ell(a)/a^\ell\le1/a\quad(\ell\ge2).
 \tag{TWP.7}
\]
The proof and the bound for all \(\ell\ge2,a\ge4\) are the
already accepted local input of HD-TW01; no new approximation to
the cycle spectrum is made here. Convolution gives at least the
weighted average \(\sum_j(\ell_j/t)s_j\).
The number of nontrivial cycles is at most \(\operatorname{ind}(\pi_i)\).
The actual natural block index is
\[
 I_i=n\operatorname{ind}(\pi_i)+\sum_j\operatorname{ind}(x_j).
 \tag{TWP.8}
\]
Using (TWP.3) only in its weaker form
\(\nu(W(x))\ge A\operatorname{ind}(x)\), valid also for \(x=1\),
therefore yields
\[
 \operatorname{rank}(1-T_i)/e
       \ge AI_i/(ta)+\lambda\operatorname{ind}(\pi_i)/t.
 \tag{TWP.9}
\]
For every factor-fixing entry the additional \(B/a\) in (TWP.5)
is available, including those that act only through the partner.

The two actual connected quotient covers satisfy Riemann--Hurwitz:
\(\sum_i I_i\ge2nt-2\) and
\(\sum_i\operatorname{ind}(\pi_i)\ge2t-2\).
Their branch sets may omit original genuine marks; the omitted
permutations contribute zero to these exact sums. Since \(A,\lambda\ge0\),
the preceding bounds and \(r_{\rm fix}\ge9\) give
\[
 \begin{split}
 \sum_i\operatorname{rank}(1-T_i)-2e
 &\ge\frac ea\left[r_{\rm fix}B-2-\frac2t(L-1)\right]\\
 &\ge\frac ea(9B-L-1)
 \ge a(9B-L-1)>6.
 \end{split}\tag{TWP.10}
\]
The second inequality uses \(L\ge1,t\ge2\). This contradicts
(TWP.1) and proves the assertion.
\end{proof}

\subsection*{HD-TW05. Two-box, fixed-tail, and cubic-trace applications}

Under (TWP.1), suppose the entire partner has support floor
\(\eta\ge1/n\). Each of the following conditions excludes the source:
\begin{enumerate}
\item \(n\ge15\) and \(W\) has label \((n-2,2)\) or
\((n-2,1,1)\), or their transposes;
\item \(n\ge24,a\ge n^3\), and
\( |\chi_W(x)|/a\le(\max\{f(x),1\}/n)^3\) for all \(x\in S_n\);
\item \(3\le k\le240\), \(W=S^{(n-k,\mu)}\), and
\(n\ge N_\mu\), with the exact accepted threshold (TW.25).
\end{enumerate}
In particular, the published effective MSP input makes the second case
apply for \(n\ge4000,a\ge n^{120}\). At the explicit integer
\(\mathcal N\) of (TW.20), every ordinary extendible label with
\(\min(n-\lambda_1,n-\lambda'_1)\ge2\) is excluded.
The last conclusion is solely for these actual permuted-core factors
with the stated entire-partner floor. The standard label is not included.

\begin{proof}
For the two-box labels use the accepted all-target local inequality
(ATX.3), whose derivation does not use a field or scalar conductor:
\[
 A=(n-5)/2,\qquad B=n/4,\qquad
 a=n(n-3)/2\ \hbox{or}\ (n-1)(n-2)/2.
 \tag{TWP.11}
\]
Here \(B/a\le1/n\) for \(n\ge6\) in the first case and
\(n^2\le2(n-1)(n-2)\) in the second, which holds for \(n\ge6\).
Also \(2B\le a-1\). The two values of \(L\) are
\((3n-5)/2\) and \((3n-3)/2\), and those of \(a\lambda\)
are \(n-1\) and \(n\). Thus
\[
 8B-2=2n-2,\qquad
 9B-L-1=(3n+6)/4\ \hbox{or}\ (3n+2)/4.
 \tag{TWP.12}
\]
Both required products with \(a\) exceed six already at \(n=15\),
and all factors increase thereafter. This checks every condition of
(TWP.4), proving the first case.

For the cubic-trace hypothesis, the complete local Fourier argument
of HD-TW02 gives
\[
 A=a(n-4)/n^2,\qquad B=a/n,
 \quad L=a(5n-4)/n^2,\quad a\lambda=4a/n-1.
 \tag{TWP.13}
\]
The verified inequalities (TW.16) give \(L\ge1,\lambda\ge0\),
\(2B\le a-1\), and \(a(9B-L-1)>6\).
Moreover \(a(8B-2)\ge a(8n^2-2)>6\).
The entire-partner floor is exactly at least \(B/a\), so HD-TW04
applies. The stated effective MSP consequence is the accepted published
input with its exact locators and numerical-reading boundary in HD-TW02.

For the fixed-tail case retain the local constants (TW.26):
\(S=Cn^{k-1}\ge100\), \(.99nS\le a\le1.01nS\),
\(A=a/n-1.01S/n\), \(B=.4S\).
Then \(B/a\le .4/(.99n)<1/n\). All the other conditions of
(TWP.4), apart from its first product, were checked explicitly in
(TW.26)--(TW.27); the first product satisfies
\(a(8B-2)=a(3.2S-2)>6\). This proves the third case.

Finally \(\mathcal N\) exceeds every \(N_\mu\), as proved in
HD-TW03, and it also exceeds \(241^{482}\): indeed
\(F_*>241^2\), hence \(\mathcal N>241^{960}\).
The branching and subpartition argument (FT.16)--(FT.17) is a pure
dimension bound, independent of any cover. It says that a label with
both tails at least241 has \(a\ge n^{120}\) at this threshold.
Such a label is excluded by the cubic-trace case. A tail between3
and240 is excluded by the fixed-tail case, and tail2 by the first case.
Transposition changes only a sign character and does not change any
projective support. These exhaust the asserted labels.
\end{proof}

\paragraph{Retained scope and significance.}
The support floor may later be supplied by an actual complete
partner decomposition, but it is an explicit hypothesis here.
No individual \(A_n\) quotient is extracted from permuted components.
Only the two genuine quotient actions on \(t\) and \(nt\) letters
are used. An unbounded omitted multiplicity space would invalidate
(TWP.2), and is not silently removed. Nonextendible alternating
constituents, spin modules, and the standard label remain outside
HD-TW05 as written. These are necessary supplier exclusions and do
not construct cycles or assert a new integral detector bound.

The result narrows the possible curve suppliers for condition A
(cote16, scan050; SIGN\_NEXT SN-C03/C04). Condition I (scan049;
cote44, Theorem4), the Cayley--Hamilton projector (scans048--049),
and the existing integral spreading interfaces retain their earlier
roles. There is no new assertion about motivic L-values, and no use
of Markman, Severn, or032. The weight-one local system calculation
is the exact Chevalley--Weil moving-rank budget of an actual cover.
